\chapter{The contest machine and the cost model}\label{ch:machine}
\index{contest workflow}

Competitive programming is a game played inside a strict machine: a judge
compiles your program, runs it on hidden tests, and returns a one-word verdict
in seconds. This chapter teaches you how that machine works and how to think
inside it. By the end you will have a starter file you trust, a working sense
of how fast a given algorithm is before you write it, and a routine for going
from a problem statement to an accepted submission without flailing. Everything
else in the book assumes this chapter's machinery.

\section{How a judge sees your program}\label{sec:judge}

Forget the IDE, the debugger, the terminal prompt. To the judge\index{judge}
your program is a \emph{function}: it reads characters from standard input
(\texttt{stdin}\index{stdin}) and writes characters to standard output
(\texttt{stdout}\index{stdout}). No files, no network, no graphics, no prompts.
A program that prints \texttt{please enter n:} gets marked wrong, because the
judge compares your output byte-for-byte with the expected output.

A submission is run once per \emph{test file}. Each test file is one complete
input; your program must print the correct answer for it within the
\emph{time limit}\index{time limit} (typically 1--2 seconds) and the
\emph{memory limit}\index{memory limit} (typically 256--512\,MB). The tests are
hidden from you. They are also written by people whose job is to break
plausible-but-wrong programs, so ``it worked on my three tests'' means nothing.

After all runs, the judge reports one \emph{verdict}\index{verdict} per
submission, sometimes per test:

\begin{center}
\begin{tabular}{ll}
\toprule
\thead{Verdict} & \thead{Meaning}\\
\midrule
Accepted (AC) & every test file passed\\
Wrong Answer (WA) & your output differed from the expected output\\
Time Limit Exceeded (TLE) & a run did not finish inside the time limit\\
Runtime Error (RE) & the program crashed (segfault, bad index, \dots)\\
Memory Limit Exceeded (MLE) & the program exceeded the memory budget\\
Compile Error (CE) & the judge could not build your code\\
\bottomrule
\end{tabular}
\end{center}

Figure~\ref{fig:verdict} shows the whole pipeline; Figure~\ref{fig:statement}
annotates a typical statement so you can see where the machine's parameters
are declared.

\begin{figure}[ht]\centering
% fig01-verdict.tex -- Figure 1.1: what the judge does with a submission.
\begin{tikzpicture}[
  font=\small,
  box/.style={draw, rounded corners=1.5pt, minimum height=7.5mm,
              inner xsep=6pt, fill=graybg},
  good/.style={draw, rounded corners=1.5pt, minimum height=6.5mm,
              inner xsep=5pt, fill=goodbg},
  bad/.style={draw, rounded corners=1.5pt, minimum height=6.5mm,
              inner xsep=5pt, fill=amberbg},
  lab/.style={font=\footnotesize, text=slate},
  arr/.style={-{stealth}, thick}
]
\node[box] (sub) at (0,0) {source code};
\node[box, right=7mm of sub] (comp) {compile};
\node[box, right=7mm of comp] (run) {run on hidden tests};
\node[good, right=7mm of run, yshift=9mm] (ac) {Accepted};
\node[bad,  right=7mm of run, yshift=-9mm] (rej) {Rejected};
\draw[arr] (sub) -- (comp);
\draw[arr] (comp) -- (run);
\draw[arr] (run) -- (ac) node[midway, above, lab] {every test passed};
\draw[arr] (run) -- (rej) node[midway, below, lab] {any failure};
\draw[arr] (comp.south) |- ($(comp.south)+(0,-5mm)$) -|
  node[lab, pos=0.75, left] {does not compile} (rej.north);
\node[lab, text width=12.4cm, align=center] at ($(run.south)+(0.4,-1.9)$)
  {Rejected comes in five flavours: \textbf{WA} wrong answer,
   \textbf{TLE} too slow, \textbf{RE} crash, \textbf{MLE} too much memory,
   \textbf{CE} never built};
\end{tikzpicture}

\caption{What the judge does with your code. Notice there is only one good
outcome and five distinct ways to lose; each failure tells you where to look
first. CE means the file did not build at all --- fix that before anything
else.}
\label{fig:verdict}
\end{figure}

\begin{figure}[ht]\centering
% fig02-statement.tex -- Figure 1.2: anatomy of a problem statement.
% The statement itself is invented for this book; the layout follows
% the CSES / USACO convention of name, limits, task, I/O, example.
\begin{tikzpicture}[font=\small]
\node[draw=ink, fill=white, text width=7.8cm, inner sep=8pt, font=\footnotesize,
      align=flush left, anchor=north west] (stmt) at (0,0) {%
  {\sansfont\bfseries Apple Sums}\quad{\color{slate}time 1.0 s \quad memory 256 MB}\\[4pt]
  \textbf{Task.} A shop sells $n$ kinds of apples. Kind $i$ costs $a_i$ rupiah.
  Find the cheapest way to buy exactly $k$ apples (any kinds, repeats allowed),
  or report $-1$ if impossible.\\[4pt]
  \textbf{Input.} First line: $n$ and $k$. Second line: $a_1,\dots,a_n$.\\[4pt]
  \textbf{Output.} One integer: the minimum total cost, or $-1$.\\[4pt]
  \textbf{Constraints.} $1 \le n \le 10^5$, $1 \le k \le 10^9$,
  $1 \le a_i \le 10^9$.\\[4pt]
  \textbf{Example.} Input \texttt{3 5 / 4 2 9} $\to$ output \texttt{10}.};
\begin{scope}[font=\footnotesize, text=accent, text width=3.7cm, align=flush left, anchor=west]
\node (c1) at (8.9,-0.45) {\textbf{limits:} the budget your program lives under};
\node (c2) at (8.9,-1.75) {\textbf{I/O contract:} what stdin gives, what stdout must print};
\node (c3) at (8.9,-3.05) {\textbf{constraints:} these numbers choose the algorithm};
\node (c4) at (8.9,-4.2) {\textbf{example:} a handshake test, not a proof};
\end{scope}
\draw[-{stealth}, accent, thick] (c1.west) -- (8.35,-0.45);
\draw[-{stealth}, accent, thick] (c2.west) -- (8.35,-1.45);
\draw[-{stealth}, accent, thick] (c3.west) -- (8.35,-2.5);
\draw[-{stealth}, accent, thick] (c4.west) -- (8.35,-3.2);
\end{tikzpicture}

\caption{Anatomy of a problem statement (invented here, in the layout CSES and
USACO use). Notice the constraints block: it is the most important paragraph
on the page, and Section~\ref{sec:constraints} teaches you to read it before
reading anything else.}
\label{fig:statement}
\end{figure}

Scoring varies by contest. CSES and most archives are all-or-nothing: a
submission is worth points only if every test passes. USACO gives partial
credit per test file. IOI-style problems and some Codeforces subtasks give
partial credit for easier constraint ranges --- you can AC the small
\emph{subtasks}\index{subtask} while the hard ones stay red. Always read which
model a problem uses; it changes which submissions are worth your time.

\begin{exercise}[easy]\label{exr:verdicts}
A friend submits three programs to the same problem. Predict each verdict:
(a) the program prints \texttt{answer: 42} instead of \texttt{42};
(b) the program is correct but takes 30 seconds on the largest test;
(c) the program declares \texttt{int a[1000000000];} inside \texttt{main}
(a four-gigabyte stack array).
\end{exercise}

\begin{example}[WE 1.1 --- CSES \emph{Weird Algorithm}: the full loop]
\label{ex:weird}
A problem asks: given an integer $n$ up to $10^6$, print the sequence where
each step halves an even number and sends an odd number to $3n+1$, ending at
$1$. Input \texttt{3} should print \texttt{3 10 5 16 8 4 2 1}.
\par\smallskip
\textbf{Strategy.} There is nothing to optimise: the statement \emph{is} the
algorithm, so simulate it directly. The only question is the integer type ---
intermediate values can exceed $10^6$ even when the input does not. (For the
record: starting from $n=704\,511$, legal under the limit, the sequence peaks
at $56\,991\,483\,520$.) A 32-bit \texttt{int} would overflow on a fraction of
tests, so use 64 bits.
\par\smallskip
\begin{solution}
\begin{lstlisting}
#include <iostream>
using namespace std;
using ll = long long;
int main() {
    ll n;
    cin >> n;
    while (n != 1) {
        cout << n << ' ';
        n = (n % 2 == 0) ? n / 2 : 3 * n + 1;
    }
    cout << 1 << '\n';
}
\end{lstlisting}
Run on the sample:
\begin{lstlisting}[numbers=none]
$ echo 3 | ./weird
3 10 5 16 8 4 2 1
\end{lstlisting}
It matches, so we also check the edge case $n=1$ (prints \texttt{1}
immediately) before submitting. The submission returns AC. The real workflow
here is what matters: read, simulate, test on the sample, test one edge case,
submit. You will repeat it hundreds of times.
\end{solution}
\end{example}

\section{The C++17 starter kit}\label{sec:starter}
\index{starter file|see{template}}
\index{template}

Every solution in this book is C++17. Listing~\ref{lst:template} is the
starter file; it lives in \texttt{code/lib/} and every later chapter
copies it.

\lstinputlisting[caption={\texttt{code/lib/template.cpp}: the starter file.
Minimal on purpose --- every line earns its place.},label={lst:template}]
{code/lib/template.cpp}

Three choices in it deserve a defence. First, the includes are explicit.
You will see \texttt{\#include <bits/stdc++.h>} in other people's code; judges
accept it, but it is a GCC-specific catch-all, it hides which facilities you
actually use, and it fails on compilers that do not ship it. We name the four
headers this chapter needs and add others as the book demands them.

Second, \texttt{using ll = long long}. The default integer type for counting
should be 64-bit: \texttt{int}\index{int@{\texttt{int}}} holds only
about $\pm 2.1\cdot 10^9$, while \texttt{ll} holds about
$\pm 9.2\cdot 10^{18}$. The price is a few characters and roughly twice the
memory per number --- cheap insurance against the single most common WA cause
in the archive.

Third, the two macros. \texttt{rep(i,a,b)} is a counted loop;
\texttt{all(x)} expands to \texttt{begin(x), end(x)} for calls like
\texttt{sort(all(v))}. We keep the macro count at two because macros are
invisible at read time --- a template that does magic for you will do magic
\emph{to} you in a later chapter's debugging session.

\begin{definition}[Integer overflow]\index{overflow}
\emph{Overflow} is arithmetic leaving the range of its type. In C++ signed
overflow is \emph{undefined behaviour}\index{undefined behaviour}: the standard
allows anything. In practice the value usually wraps to garbage
($2\,147\,483\,647+1$ becomes $-2\,147\,483\,648$), but ``usually'' is the
point --- the same code can misbehave differently across compilers and flags.
\end{definition}

The rules of thumb: any product of two inputs near $10^5$ or larger, any sum
of up to $10^5$ terms near $10^9$, and any expression like \texttt{n*(n+1)/2}
with $n$ above $46\,000$ exceeds \texttt{int}. When in doubt, \texttt{ll}.

\begin{tripwire}[The overflow tripwire]
If a program passes the samples but fails the judge's big tests, check the
types of every accumulator and every product \emph{before} anything else.
Worked Example 1.4 is this exact failure.
\end{tripwire}

\subsection{Fast input and output}\label{sec:fastio}
\index{fast input/output}
\index{ios::sync\_with\_stdio@\texttt{ios::sync\_with\_stdio}}

By default, C++ streams are synchronised with the C stdio library --- a
convenience you pay for on every \texttt{cin}. Two lines in the template turn
that off: \texttt{ios::sync\_with\_stdio(false)} detaches the streams, and
\texttt{cin.tie(nullptr)} stops \texttt{cin} from flushing \texttt{cout}
before each read. Together they make stream I/O fast enough for almost every
problem in this book.

When input gets large --- a million tokens or more --- even that can be the
bottleneck. Listing~\ref{lst:fastscan} compares the two levels on a file of
one million integers: plain \texttt{cin} with sync disabled versus a small
scanner that reads the raw \texttt{stdin} buffer itself
(\texttt{code/lib/fastscan.h}, part of the capstone library).

\lstinputlisting[caption={\texttt{code/fastscan.cpp}: two ways to read a
million integers, timed.},label={lst:fastscan}]{code/fastscan.cpp}

\begin{checkpoint}[Predict before you run]
Predict the ratio of the two timings: a factor of 2, of 10, or of 100? Then
look at the measured run below.
\end{checkpoint}

\begin{lstlisting}[numbers=none]
$ ./fastscan < ints1e6.txt
method=cin sum=499884612927859 time_ms=919.0
$ ./fastscan f < ints1e6.txt
method=fread sum=499884612927859 time_ms=6.6
\end{lstlisting}

On this machine the buffer scanner is about $140\times$ faster --- nine
hundred milliseconds spent purely on parsing. Take the lesson as a range, not
a law: the ratio varies by machine and judge, but the direction never does.
Use the template's two lines always; reach for \texttt{fastscan.h} when input
is the bottleneck.

\subsection{The STL you actually use}\label{sec:stl}
\index{STL}\index{vector@\texttt{vector}}\index{sort@\texttt{sort}}

You can solve the first hundred problems of this book with five library
tools. \texttt{vector} is the only array you need: contiguous, resizable,
O(1) index access. \texttt{sort} orders a range in $O(n\log_2 n)$ and accepts
a \emph{comparator}\index{comparator} --- a function returning true when the
first argument belongs before the second. \texttt{set}\index{set@\texttt{set}}
keeps distinct elements sorted (insert, erase, find all $O(\log_2 n)$);
\texttt{map}\index{map@\texttt{map}} is the same idea keyed by a value, which
makes it a frequency table; \texttt{priority\_queue}\index{priority queue}
keeps the current maximum on top in $O(\log_2 n)$ per operation. Strings are
\texttt{std::string}; small records are \texttt{pair} or \texttt{tuple},
unpacked by a \emph{structured binding}\index{structured binding}
(\texttt{auto\& [k, v] : m}). Listing~\ref{lst:stltour} exercises all of them.

\lstinputlisting[caption={\texttt{code/stl\_tour.cpp}: the four containers
and a custom comparator, with real output.},label={lst:stltour}]
{code/stl_tour.cpp}

\begin{lstlisting}[numbers=none]
sorted: 2 4 8 1 5
set: 1 2 4
a=2 b=3 c=1 
pop order: 9 7 4 3
\end{lstlisting}

Notice the comparator's effect on the first line: evens sorted before odds,
each group ascending. Notice too that \texttt{freq[w]++} needs no
initialisation --- \texttt{map} value-initialises missing entries to zero,
which is why frequency tables are one line of code.

\begin{exercise}[easy]\label{exr:predict}
Predict the output of this program, then check yourself by running it:
\begin{lstlisting}
int a = 7;
cout << a/2 << ' ' << a/2.0 << ' ' << (a*a)/2 << '\n';
\end{lstlisting}
What type does the middle expression evaluate in, and why do the first and
last answers differ?
\end{exercise}

\section{Time complexity}\label{sec:complexity}
\index{time complexity}

``Is this algorithm fast enough?'' has to be answerable \emph{before} you
write code, or you will spend contest minutes implementing something that was
doomed from the statement. The tool is asymptotic analysis: count how the
work grows with the input size, not how many milliseconds it takes on one
machine.

\begin{definition}[RAM model]\index{RAM model}
We count \emph{operations}: arithmetic, comparisons, assignments, array
reads and writes each cost one unit. A function call is the sum of the work
it does. This is a model, not physics --- real machines have caches,
branches and vector units --- but it predicts the shape of runtimes far
better than it predicts their absolute values.
\end{definition}

Count the operations in a nested loop exactly once and you will see why the
abstraction helps:

\begin{lstlisting}
long long s = 0;                       // 1 op
for (int i = 0; i < n; ++i)            // n iterations
    for (int j = i + 1; j < n; ++j)    // n-i-1 iterations
        s += a[i] * a[j];              // 3 ops each pass
\end{lstlisting}

The innermost line runs $\sum_{i=0}^{n-1}(n-1-i) = n(n-1)/2$ times, so the
total is $1 + 3n(n-1)/2$. For large $n$ the $n^2$ term dominates: at
$n=10^5$ the count is about $1.5\cdot 10^{10}$ --- seconds to minutes, not
the microseconds a linear algorithm needs. The lower-order terms and the
constants stop mattering quickly, which is exactly what the notation
formalises.

\begin{definition}[$O$, $\Theta$, $\Omega$]\index{big-O}\index{Theta@$\Theta$}%
\index{Omega@$\Omega$}
For functions $f,g:\mathbb{N}\to\mathbb{R}_{>0}$:
\begin{itemize}
\item $f \in O(g)$: there exist $c>0$ and $n_0$ with
$f(n) \le c\,g(n)$ for all $n\ge n_0$ --- $g$ is an \emph{asymptotic upper
bound}.
\item $f \in \Omega(g)$: there exist $c>0$ and $n_0$ with
$f(n) \ge c\,g(n)$ for all $n\ge n_0$ --- an asymptotic lower bound.
\item $f \in \Theta(g)$: both hold --- $f$ and $g$ have the same growth rate.
\end{itemize}
Writing $f = O(g)$ is common shorthand for $f \in O(g)$.
\end{definition}

\begin{theorem}\label{thm:thetapoly}
$3n^2 + 10n + 7 \in \Theta(n^2)$.
\end{theorem}

\begin{proof}
Upper bound: for $n \ge 1$, $10n \le 10n^2$ and $7 \le 7n^2$, so
$3n^2+10n+7 \le 20n^2$; take $c=20$, $n_0=1$ and $f \in O(n^2)$.
Lower bound: for $n \ge 1$, $3n^2+10n+7 \ge 3n^2$; take $c=3$, $n_0=1$
and $f \in \Omega(n^2)$. Both directions hold, so $f \in \Theta(n^2)$.
\end{proof}

The same squeeze works for any polynomial: the leading term is the whole
story. Inside $O(\cdot)$ the base of a logarithm never matters either, since
$\log_a n = \log_b n / \log_b a$ differs only by a constant; in this book
complexity uses $\log_2$ to keep the concrete meaning --- the number of
halvings --- and every other $\log$ is natural.

The classes that matter, from cheapest to most expensive:

\begin{center}
\begin{tabularx}{\linewidth}{@{}lYY@{}}
\toprule
\thead{Class} & \thead{Name} & \thead{Produced by}\\
\midrule
$O(1)$ & constant & a formula, one array lookup\\
$O(\log_2 n)$ & logarithmic & halving the input each step (binary search)\\
$O(\sqrt{n})$ & square root & per-block passes in sqrt decomposition\\
$O(n)$ & linear & one pass over the input\\
$O(n\log_2 n)$ & log-linear & sorting; divide and conquer without blowup\\
$O(n^2)$ & quadratic & all pairs; doubly nested loops\\
$O(n^3)$ & cubic & all triples; some DP on small $n$\\
$O(2^n)$ & exponential & enumerating all subsets\\
$O(n!)$ & factorial & enumerating all permutations\\
\bottomrule
\end{tabularx}
\end{center}

\begin{figure}[ht]\centering
% fig03-growth.tex — Figure 1.3: operation-count growth vs the 1 s budget.
% Log-log axes; the 1e8-op line marks one second of simple work.
\begin{tikzpicture}
\begin{loglogaxis}[
    width=0.86\linewidth, height=6.8cm,
    xlabel={input size $n$}, ylabel={operations},
    xmin=1, xmax=1e8, ymin=1, ymax=1e17,
    legend pos=south east, legend style={font=\footnotesize},
    grid=major, grid style={gray!22},
    tick label style={font=\footnotesize},
    label style={font=\small}]
\addplot[domain=1:1e8, samples=60, thick, accent] {ln(x)/ln(2)};
\addplot[domain=1:1e8, samples=60, thick] {x};
\addplot[domain=1:1e8, samples=80, thick, dashed] {x*ln(x)/ln(2)};
\addplot[domain=1:1e8, samples=60, thick, dotted] {x^2};
\addplot[domain=1:60,   samples=61, thick, dashdotted] {2^x};
\addplot[domain=1:1e8, samples=2, thin, amber] {1e8};
\legend{$\log_2 n$, $n$, $n\log_2 n$, $n^2$, $2^n$, 1-second budget ($10^8$)}
\end{loglogaxis}
\end{tikzpicture}

\caption{Operation counts on log--log axes against the one-second budget of
$10^8$ simple operations. Notice where each curve crosses the amber line:
$\log_2 n$ never gets there, $n\log_2 n$ crosses near $n \approx 4\cdot10^6$,
$n^2$ near $10^4$, and $2^n$ already fails before $n=30$. This single picture
is the whole point of Section~\ref{sec:constraints}.}
\label{fig:growth}
\end{figure}

\begin{definition}[The $10^8$ heuristic]\index{10^8 heuristic@$10^8$ heuristic}
As a budgeting rule of thumb, one second of judge time buys roughly $10^8$
\emph{simple} operations --- integer adds, compares, array touches. It is a
heuristic, not a law: memory-bound passes can be slower, vectorised inner
loops can be several times faster, and the number says nothing about a
two-second or ten-second limit. Use it to compare candidate algorithms, not
to certify one.
\end{definition}

\begin{exercise}[medium]\label{exr:countops}
Count the operations in this fragment exactly, then give its $\Theta$-class:
\begin{lstlisting}
for (int i = 1; i <= n; ++i)
    for (int j = 1; j <= i; j += 2)
        work();
\end{lstlisting}
Hint: the inner trip count is $\lceil i/2 \rceil$.
\end{exercise}

\section{From constraints to a plan}\label{sec:constraints}
\index{constraints}

Statements hand you the input size $n$ on a plate. The contest skill is
converting $n$ into the \emph{largest} complexity class that survives the
budget --- Figure~\ref{fig:constraints} is the table every experienced
contestant has internalised --- and then picking the simplest algorithm in
that class.

\begin{figure}[ht]\centering
% fig04-constraints.tex -- Figure 1.4: the constraint-to-complexity table,
% drawn as a boxed chart (a figure, not body text).
\renewcommand{\arraystretch}{1.25}%
\begin{tabular}{@{}r l l@{}}
\toprule
\thead{largest $n$} & \thead{needs at most} & \thead{typical shapes}\\
\midrule
$n \le 12$      & $O(n!)$                 & permutations, brute force\\
$n \le 25$      & $O(2^n)$                & subsets, meet-in-the-middle\\
$n \le 500$     & $O(n^3)$                & all triples, $n\times n$ DP\\
$n \le 5\,000$  & $O(n^2)$                & all pairs, quadratic DP\\
$n \le 2\cdot 10^5$ & $O(n\log_2 n)$    & sorting, sets, Fenwick\\
$n \le 10^7$    & $O(n)$                  & single linear pass\\
any $n$         & $O(\log_2 n)$ or $O(1)$ & binary search, closed formula\\
\bottomrule
\end{tabular}%
\\[4pt]
\begin{minipage}{0.9\linewidth}\footnotesize\color{slate}
Assumes the $10^8$-ops-per-second budget. Rows are rules of thumb, not laws:
a $5\cdot 10^7$-constant quadratic can still lose to a linear algorithm.
\end{minipage}

\caption{The constraint-to-complexity table. Notice it is a ceiling chart:
$n\le 2\cdot10^5$ does not demand $O(n\log_2 n)$, it \emph{forbids} anything
slower. Rows are rules of thumb under the $10^8$ heuristic; a
$5\cdot10^7$-constant quadratic can still lose to a linear pass.}
\label{fig:constraints}
\end{figure}

Two refinements before the table gets used. First, judge tests are built to
stress the \emph{worst case}\index{worst case}: average-case reasoning is how
quicksort dies on sorted input. Plan for the worst row in Figure's table,
not the typical one. Second, \emph{constant factors}\index{constant factor}
are real --- an $O(n\log_2 n)$ algorithm whose inner step rebuilds a
\texttt{set} can lose to a well-written $O(n)$ pass at the same $n$ --- but
constants lose to growth: no constant rescues $n^2$ when $n=10^6$.

\begin{rulebox}[Choosing the algorithm]
Let $C$ be the largest class that fits the constraint table for the given
$n$. Among algorithms in $C$, implement the \emph{simplest correct one you
can explain} --- the one with the fewest moving parts. If it TLEs, only then
reach for the cleverer algorithm in the same class. Never climb the table
first: a fancy $O(n\log_2 n)$ where $O(n^2)$ would have passed is wasted
minutes.
\end{rulebox}

\begin{example}[WE 1.2 --- CSES \emph{Maximum Subarray Sum}, four ways]
\label{ex:mss}
Given an array of $n$ integers (with negatives), find the largest sum of a
contiguous subarray. On CSES, $n \le 2\cdot10^5$: the table says
$O(n\log_2 n)$ or better.
\par\smallskip
\textbf{Strategy.} Same problem, four honest attempts, each removing one
structural waste: (i) try every interval and resum --- correct but cubic;
(ii) keep the sum running while the right end moves --- quadratic;
(iii) split the array and combine --- $O(n\log_2 n)$;
(iv) notice that only the \emph{best ending} matters at each step ---
linear.
\par\smallskip
\begin{solution}
The cubic version checks all $n(n+1)/2$ intervals and resums each:
\begin{lstlisting}
for (int l = 0; l < n; ++l)
    for (int r = l; r < n; ++r) {
        ll s = 0;
        for (int k = l; k <= r; ++k) s += a[k];
        best = max(best, s);
    }
\end{lstlisting}
It is correct --- it inspects every candidate --- and it is the first
submission most beginners write. The quadratic version keeps $s$ running as
$r$ advances, deleting the inner loop. The divide-and-conquer version splits
at the middle: the answer is the best of left half, right half, or an
interval crossing the middle (extend left and right greedily from the
split). The linear version is Kadane's\index{Kadane's algorithm}:
\begin{lstlisting}
ll best = a[0], cur = a[0];
for (int i = 1; i < n; ++i) {
    cur = max(a[i], cur + a[i]);   // extend or restart
    best = max(best, cur);
}
\end{lstlisting}
Correctness, in one line: \texttt{cur} is the largest sum of any subarray
\emph{ending at} $i$ --- you either continue the previous best suffix or
abandon it --- and \texttt{best} remembers the largest of those. All four
are in \texttt{code/mss\_timing.cpp}; Listing~\ref{lst:msstiming} is the
driver. Measured on this machine:
\begin{lstlisting}[numbers=none]
O(n^3)    n=3000    ans=24382 time_ms=254.9
O(n^2)    n=50000   ans=109297 time_ms=602.2
O(nlogn)  n=200000  ans=218662 time_ms=1.8
O(n)      n=200000  ans=218662 time_ms=0.1
O(n)      n=5000000 ans=1565779 time_ms=3.4
\end{lstlisting}
Read the table backwards: the linear pass handles $n=5\cdot10^6$ in less
time than the cubic one needs for $n=3000$ --- and note the \texttt{ans}
column is identical wherever both finished, which is the definition of a
correct optimisation.
\end{solution}
\end{example}

\lstinputlisting[caption={Listing 1.4: the benchmark driver
(\texttt{code/mss\_timing.cpp}, excerpt). Each algorithm is timed at the
largest $n$ it survives in about a second; the linear one is pushed to
$n=5\cdot10^6$ to show headroom.},
label={lst:msstiming},firstline=63,lastline=85]{code/mss_timing.cpp}

We did not run the cubic version at $n=2\cdot10^5$; extrapolating the
measured scaling, it would take about a day. The quadratic at that $n$ is
roughly ten seconds. Both lose the contest to algorithm choice, not to
coding skill --- which is the lesson.

\begin{exercise}[medium]\label{exr:pickclass}
For each problem, name the largest complexity class the constraint allows
under the $10^8$ heuristic and the shape of an algorithm that fits:
(a) $n \le 15$: count the subsets with a given property;
(b) $n \le 3\cdot10^5$: find the maximum of a sliding window;
(c) $n \le 2000$: longest common subsequence of two strings.
\end{exercise}

\section{The submit loop}\label{sec:submit}
\index{submit loop}
\index{sample test}
\index{Wrong Answer|see{verdict}}

Acceptance is a loop, not an event. Figure~\ref{fig:minute} shows how a
working five minutes splits on one problem; the dashed arrow is the part
beginners get wrong --- a WA sends you back to \emph{thinking}, not to
resubmitting.

\begin{figure}[ht]\centering
% fig05-minute.tex — Figure 1.5: how a working five minutes is spent on
% one problem. The same loop returns in Chapter 18 (contest strategy).
\begin{tikzpicture}[font=\small,
  seg/.style={draw, minimum height=8mm, inner xsep=0pt},
  lab/.style={font=\footnotesize},
  num/.style={font=\footnotesize, text=slate}]
% proportional widths: read 30s think 75s code 100s test 45s submit 50s = 300s
\def\u{0.035} % seconds -> cm
\node[seg, fill=graybg,   minimum width=30*\u cm,  anchor=west] (a) at (0,0) {};
\node[seg, fill=accentsoft, minimum width=75*\u cm, anchor=west] (b) at (a.east) {};
\node[seg, fill=graybg,   minimum width=100*\u cm, anchor=west] (c) at (b.east) {};
\node[seg, fill=accentsoft, minimum width=45*\u cm, anchor=west] (d) at (c.east) {};
\node[seg, fill=goodbg,   minimum width=50*\u cm,  anchor=west] (e) at (d.east) {};
\node[lab] at (a.center) {read};
\node[lab] at (b.center) {think};
\node[lab] at (c.center) {code};
\node[lab] at (d.center) {test};
\node[lab] at (e.center) {submit};
\node[num, below=1pt of a.west, anchor=north west] {0:30};
\node[num, below=1pt of b.center] {1:45};
\node[num, below=1pt of c.center] {3:25};
\node[num, below=1pt of d.center] {4:10};
\node[num, below=1pt of e.east, anchor=north east] {5:00};
\draw[{stealth}-, thick, accent] (c.north) -- ++(0,0.5)
  node[above, font=\footnotesize, text=accent]
  {longest block is still under two minutes --- pick the simple plan};
\draw[-{stealth}, dashed, slate] (e.south east) ++(0,-1.1)
  -| node[pos=0.25, below, font=\footnotesize, text=slate]
  {WA $\to$ back to ``think'', not to ``submit again''} (b.south);
\end{tikzpicture}

\caption{One problem, five minutes. Notice the blocks are small and the
feedback edge goes to ``think'': a wrong answer is new information about
your assumptions, not an invitation to submit again.}
\label{fig:minute}
\end{figure}

The loop, expanded. \emph{Compile like the judge}: the book's flags are
\texttt{-std=c++17 -O2 -Wall -Wextra}; warnings are unheeded verdicts.
\emph{Test locally first}: run the samples, then one edge case you invented
($n=1$, all-equal input, the largest legal $n$ if generation is cheap).
\emph{Read a WA as a clue}, in this order:

\begin{checklist}
\item Did I handle the constraints' endpoints ($n=1$, $n=0$, maximum $n$)?
\item Are all accumulators and products 64-bit? (The overflow tripwire.)
\item Does my output match the format byte-for-byte --- no extra words,
spaces, or missing newline?
\item Did I test anything the statement did not give me, or did I trust
the sample?
\end{checklist}

And keep one thing straight: \textbf{Accepted means passed these tests, not
correct}. Judges write strong test sets, but ``AC'' certifies nothing about
inputs nobody tried. The deep fix --- proving your candidate against a brute
force on random small tests --- is \emph{stress testing}\index{stress test},
and Chapter~18 makes it a habit; \texttt{code/test\_ch01.sh} already runs
200 random rounds of the linear MSS code against the cubic one.

\begin{example}[WE 1.3 --- CSES \emph{Missing Number}: two ways to be right]
\label{ex:missing}
Numbers $1$ to $n$ appear once each except one, which is missing. Given
$n \le 2\cdot10^5$ and the $n-1$ present numbers, find the missing one.
\par\smallskip
\textbf{Strategy.} Two routes present themselves: a \emph{formula} ---
the total should be $n(n+1)/2$, so subtract the given sum --- and a
\emph{marking} pass --- record which values appear and find the gap. Both
are $O(n)$, both fit easily. Under time pressure choose the one with fewer
ways to be wrong: the formula's product can overflow if types slip, while
marking has no arithmetic at all beyond indexing. Marking is the safer
submission; the formula is the faster insight. Knowing both lets you pick.
\par\smallskip
\begin{solution}
Formula (\texttt{code/missing\_math.cpp}) and marking
(\texttt{code/missing\_set.cpp}); on the test of Section~1.5 both print
\texttt{200000} for the crafted $n=2\cdot10^5$ case. The point is not the
answer --- it is that two correct algorithms existed and the choice was a
risk decision, not a knowledge gap.
\end{solution}
\end{example}

\begin{example}[WE 1.4 --- the first bug hunt]
\label{ex:bughunt}
Here is a real Wrong Answer. The sum formula for \emph{Missing Number}
written with \texttt{int} throughout (\texttt{code/missing\_int.cpp}) passes
the sample happily --- and prints this on the crafted large test:
\begin{lstlisting}[numbers=none]
$ ./missing_int < missing_big.txt
-2147283648
\end{lstlisting}
Negative two billion is not a number between $1$ and $n$.
\par\smallskip
\textbf{Strategy.} Apply the checklist in order. Endpoints: $n$ legal, fine.
Types: \texttt{int total = n*(n+1)/2} overflows once
the intermediate product $n(n+1) > 2^{31}$, i.e.\ already at
$n \approx 46\,341$ --- the judge's tests go to $2\cdot10^5$.
That is the bug, found in one checklist step.
\par\smallskip
\begin{solution}
Change three \texttt{int}s to \texttt{ll} (\texttt{code/missing\_math.cpp});
the same test now prints \texttt{200000}. The sample never had a chance to
catch this: with $n=4$, $n(n+1)/2 = 10$ fits in 16 bits, let alone 32.
``Samples pass'' and ``types are right'' are independent facts --- only the
checklist connects them.
\end{solution}
\end{example}

\begin{exercise}[medium]\label{exr:fixwa}
This program prints a million lines and gets TLE although it does almost no
work per line. Find the one-token fix and say why it matters:
\begin{lstlisting}
for (int i = 0; i < 1000000; ++i)
    cout << i << endl;
\end{lstlisting}
\end{exercise}

\section{Common mistakes}
\index{common mistakes}

\begin{itemize}
\item \textbf{\texttt{endl} in a loop.} \texttt{endl} writes a newline
\emph{and flushes} the stream; a flush per line turns microseconds into
seconds. Use \texttt{'\textbackslash n'}, and let \texttt{cin.tie(nullptr)}
handle the rest.
\item \textbf{\texttt{int} where the math needs 64 bits.}
\texttt{n*(n+1)/2} overflows at $n \approx 46\,000$; products of two
$10^5$-scale inputs overflow always. When an accumulator can reach $10^9$,
it is \texttt{ll}. See WE~1.4.
\item \textbf{Reading before understanding the format.} A common failure is
slurping tokens before noticing the input is a grid, or that the first
number is a \emph{test count} and the loop is per-case. Sketch the parse on
paper before coding.
\item \textbf{Believing the sample proves the program.} A sample is a
handshake test, not a proof: WE~1.4's bug was invisible at $n=4$. The
checklist exists precisely because samples cannot see overflows, edge
cases, or the second half of a compound condition.
\end{itemize}

\section{Chapter summary}
\begin{table}[ht]\centering
\begin{tabularx}{\linewidth}{@{}lYY@{}}
\toprule
\tablehead\thead{Item} & \thead{Rule of thumb} & \thead{First check when wrong}\\
\midrule
Verdicts & AC = all tests; WA = wrong bytes; TLE = too slow; RE = crash & Read verdict before rereading code\\
Types & $\texttt{int} < 2.1\cdot10^9$; \texttt{ll} $< 9.2\cdot10^{18}$ & Products and accumulators first\\
I/O & sync off + \texttt{tie(nullptr)} always; raw scanner past $10^6$ tokens & \texttt{endl} $\to$ \texttt{'\textbackslash n'}\\
Complexity & $10^8$ simple ops $\approx$ 1 s & count the innermost loop\\
Planning & $n$ $\to$ largest class in Fig.~\ref{fig:constraints} $\to$ simplest algorithm in it & worst case, never average\\
Submit loop & sample $\to$ edge case $\to$ submit $\to$ checklist & WA goes back to ``think''\\
\bottomrule
\end{tabularx}
\caption{Chapter 1 in one page. Keep this open during your first ten
submissions; afterwards it is muscle memory.}
\label{tab:ch1summary}
\end{table}

\section{Additional exercises}
\index{exercises}
Named archive problems are cited by name only; look them up on
\texttt{cses.fi} and \texttt{usaco.org}.

\begin{exercise}[easy] CSES \emph{Repetitions}: longest run of equal
characters. Trains a single linear pass. [implement]
\end{exercise}
\begin{exercise}[easy] CSES \emph{Increasing Array}: minimum increments to
make the array non-decreasing. Trains a one-pass greedy.
\end{exercise}
\begin{exercise}[easy] CSES \emph{Bit Strings}: count binary strings of
length $n$, answer mod $10^9+7$. First taste of modular arithmetic.
\end{exercise}
\begin{exercise}[easy] CSES \emph{Number Spiral}: value at $(y,x)$ in the
spiral. [derive a closed form --- no loops needed]
\end{exercise}
\begin{exercise}[easy] CSES \emph{Palindrome Reorder}: reorder into a
palindrome or report impossible. Trains parity counting and a decision.
\end{exercise}
\begin{exercise}[easy] USACO Bronze \emph{The Cow-Signal}: scale a grid by a
factor. Trains I/O loops and indexing.
\end{exercise}
\begin{exercise}[easy] USACO Bronze \emph{Speeding Ticket}: compare two
piecewise-constant speeds over a road. Trains interval bookkeeping.
\end{exercise}
\begin{exercise}[medium] CSES \emph{Trailing Zeros}: trailing zeros of
$n!$. [derive the formula --- then prove it]
\end{exercise}
\begin{exercise}[medium] CSES \emph{Coin Piles}: can two piles be emptied
by moves $(-2a,-a)$? [derive the test on paper first]
\end{exercise}
\begin{exercise}[medium] USACO Bronze \emph{The Bucket List}: max cows
milking at once from interval schedules. Trains sweep-line counting
(preview of Chapter~5).
\end{exercise}
\begin{exercise}[medium] USACO Bronze \emph{The Bovine Shuffle}: undo a
position permutation. Trains array-as-function thinking.
\end{exercise}
\begin{exercise}[hard] CSES \emph{Two Sets}: partition $1..n$ into two
equal-sum sets or report impossible. [construct the split --- constructive
proof required]
\end{exercise}

\chapternotes
The judge model and the submit-loop discipline follow the contest-craft
chapters of Sannemo~\cite{sannemo2025} and Laaksonen~\cite{laaksonen2024};
both books open with the same machine-first framing. The complexity
treatment --- RAM model, the squeeze proof for $\Theta$, and the
constraint table --- parallels Laaksonen ch.~2~\cite{laaksonen2024} and the
USACO Guide's Bronze modules on time complexity and
simulation~\cite{usacoguide}. The starter-file and fast-I/O conventions are
the archive's de facto standard, documented at~\cite{cppreference} and in
the CSES problem set's own submission rules~\cite{cses}; the
\texttt{fastscan} routine is the standard contest idiom also catalogued at
cp-algorithms~\cite{cpalgorithms}. Halim's text~\cite{halimcp4} covers the
same ground with a different exercise mix --- cited here, not consulted for
content. All archive problems named in this chapter come from the CSES
introductory set and the USACO Bronze archive; the rating-band mapping used
throughout the book was computed from the public Codeforces problemset and
is documented in \texttt{docs/sources.md}.
